% !TeX root = elegantnote-en.tex
\section{Carboxylic Acids}

\subsection{Introduction 20260512}

\begin{enumerate}
    \item The combination of a carbonyl group and a hydroxyl on the same carbon atom is called a \textbf{carboxyl group}.
    \item An \textbf{aliphatic acid} has an alkyl group bonded to the carboxyl group, and \textbf{an aromatic} acid has an aryl group. \textbf{Fatty acids} are long-chain aliphatic acids derived from the hydrolysis of fats and oils.
    \item A carboxylic acid donates protons by heterolytic cleavage of the acidic \chemfig{O-H} bond to give a proton and a \textbf{carboxylate ion}.
\end{enumerate}

\subsection{Nomenclature of Carboxylic Acids 20260512}

\subsubsection{Commonn Names}

\begin{enumerate}
    \item Few common names: % 12~18
\end{enumerate}

\subsubsection{IUPAC Names}

\begin{enumerate}
    \item The IUPAC nomenclature for carboxylic acids uses the name of the alkane/alkene that corresponds to the longest continuous chain of carbon atoms. The final \textit{-e} in the alkane name is replaced by the suffix \textit{-oic acid}.
    \item Cycloalkanes with \ce{-COOH} substituents are generally named as \textit{cycloalkanecarboxylic acids}. e.g. 3,3-dimethylcyclohexanecarboxylic acid
    \item Aromatic acids of the form \chemfig{Ar-COOH} are named as derivatives of \textit{benzoic acid}, \chemfig{Ph-COOH}.
\end{enumerate}

\subsubsection{Nomenclature of Dicarboxylic Acids}

\begin{enumerate}
    \item Common Names of Dicarboxylic Acids
    \begin{enumerate}
        \item A dicarboxylic acid (also called a \textit{diacid}) is a compound with two \textbf{carboxyl groups}.
        \item oxalic, malonic, maleic (can lose water), fumaric (cannot lose water), phthalic
        \begin{table}[h]
            \centering
            \begin{tabular}{|l|l|}
            \hline
                1 & 2 \\ \hline
        
            \end{tabular}
        \end{table}
    \end{enumerate} 
    \item IUPAC Names of Dicarboxylic Acids
    \begin{enumerate}
        \item Aliphatic dicarboxylic acids are named simply by adding the suffix \textit{-dioic acid} to the name of the parent alkane.
    \end{enumerate}
\end{enumerate}

\begin{example} %20-1
    Draw the structures of the following carboxylic acids.
\end{example}

\begin{example}
    Name the following carboxylic acids (when possible, give both a common name and a systematic name).
\end{example}

\subsection{Structure and Physical Properties of Carboxylic Acids 20260512}

\begin{enumerate}
    \item Structure of the Carboxyl Group
    \begin{enumerate}
        \item The entire molecule is approximately planar. The \textit{sp}$^2$ hybrid carbonyl carbon atom is planar, with nearly trigonal bond angles. The \chemfig{O-H} bond also lies in this plane, eclipsed with the \chemfig{C=O} bond.
        \item One of the unshared electron pairs on the hydroxyl oxygen atom is delocalized into the electrophilic pi system of the carbonyl group.
    \end{enumerate}
    \item Boiling Points
    \begin{enumerate}
        \item Carboxylic acids boil at considerably higher temperatures than do alcohols, ketones, or aldehydes of similar molecular weights.
        \item The high boiling points of carboxylic acids result from formation of a stable, hydrogenbonded dimer.
    \end{enumerate}
    \item Melting Points
    \item Solubilities
\end{enumerate}

\subsection{Acidity of Carboxylic Acids 20260512}

\subsubsection{Measurement of Acidity}

\begin{enumerate}
    \item A carboxylic acid may dissociate in water to give a proton and a carboxylate ion. The equilibrium constant $K_{\mathrm{a}}$ for this reaction is called the \textit{acid-dissociation constant}.
    \item benzoic acid $\mathrm{p}K_\mathrm{a} = 4.19$ phenol $\mathrm{p}K_\mathrm{a} \sim 10$ 
\end{enumerate}

\subsubsection{Substituent Effects on Acidity}

\begin{enumerate}
    \item The magnitude of a substituent effect depends on its distance from the carboxyl group.
    \item Substituted benzoic acids show similar trends in acidity, with electron-withdrawing groups enhancing the acid strength and electron-donating groups decreasing the acid strength. These effects are strongest for substituents in the ortho and para positions.
\end{enumerate}

\begin{example} % 20-3
    Rank the compounds in each set in order of increasing acid strength.
    \begin{enumerate}[label=(\alph*)]
        \item 1
        \item 2
        \item 3 \ce{-NO2} group strongest
    \end{enumerate}
\end{example}

\subsection{Salts of Carboxylic Acids 20260512}

\begin{enumerate}
    \item A strong base can completely deprotonate a carboxylic acid. The combination of a carboxylate ion and a cation is a \textbf{salt of a carboxylic acid}.
    \item Nomenclature of Carboxylic Acid Salts
    \begin{enumerate}
        \item Salts of carboxylic acids are named simply by naming the cation, then naming the carboxylate ion by replacing the \textit{-ic acid} part of the acid name with \textit{-ate}.
    \end{enumerate}
    \item Properties of Acid Salts
\end{enumerate}

\subsection{Commercial  Sources of Carboxylic Acids 20260512}

\begin{enumerate}
    \item Hydrolysis of a fat or an oil gives a mixture of the salts of straight-chain fatty acids.
    \item Important Acids
    \begin{enumerate}
        \item The most important commercial aliphatic acid is acetic acid.
        \item Aromatic acids
        \item Two important commercial diacids are adipic acid (hexanedioic acid) and terephthalic acid (benzene-1,4-dicarboxylic acid).
    \end{enumerate}
\end{enumerate}

\subsection{}
\subsection{}
\subsection{}
% watch the recording
\subsection{Condensation of Acids with Alcohols: The Fischer Esterification 20260514}

\begin{enumerate}
    \item The \textbf{Fischer esterification} converts carboxylic acids and alcohols directly to esters by an acid-catalyzed nucleophilic acyl substitutison.
\end{enumerate}

\subsection{Esterification Using Diazomethane 20260519}

\begin{enumerate}
    \item Carboxylic acids are converted to their methyl esters very simply by adding an ether solution of diazomethane.
    \item Diazomethane is a toxic, explosive yellow gas that dissolves in ether and is fairly safe to use in ether solutions.
\end{enumerate}

\begin{mechanism}
    Esterification Using Diazomethane
\end{mechanism}

\subsection{Condensation of Acids with Amines: Direct Synthesis of Amides 20260519}

\begin{enumerate}
    \item The initial acid–base reaction of a carboxylic acid with an amine gives an ammonium carboxylate salt.
    \item Heating this salt to well above \qty{100}{\degreeCelsius} drives off steam and forms an amide.
\end{enumerate}

\subsection{Reduction of Carboxylic Acids 20260519}

\begin{enumerate}
    \item Lithium aluminum hydride (\ce{LiAlH4} or LAH) reduces carboxylic acids to primary alcohols.
    \item Possible mechanism
    \item Borane also reduces carboxylic acids to primary alcohols.
    \item React with the carboxyl group faster than with any other carbonyl function, elcellent selectivity.
    \item Reduction to Aldehydes
    \begin{enumerate}
        \item Lithium tri-tert-butoxyaluminum hydride, \ce{LiAlH(O-\textit{t}-Bu)3} is a weaker reducing agent than lithium aluminum hydride. 
        \item Reduction of an acid to an aldehyde is a two-step process: Convert the acid to the acid chloride, then reduce using lithium tri-tert-butoxyaluminum hydride.
    \end{enumerate}
\end{enumerate}

\subsection{Alkylation of Carboxylic Acids to Form Ketones 20260519}

\begin{enumerate}
    \item Carboxylic acids react with two equivalents of an organolithium reagent to give ketones.
    \item Girgnard reagents cannot be used. why?
    \item Mechanism
\end{enumerate}

\subsection{Synthesis and Use of Acid Chlorides 20260519}

\begin{enumerate}
    \item Halide ions are excellent leaving groups for nucleophilic acyl substitution. Acyl halides are useful intermediates for making acid derivatives.
    \item Acid chlorides react with a wide range of nucleophiles.
    \item The best reagents for converting carboxylic acids to acid chlorides are thionyl chloride (\ce{SOCl2}) and oxalyl chloride [\ce{(COCl)2}].
    \item Mechanism
    \item Acid chlorides react with alcohols to give esters through a nucleophilic acyl substitution by the addition–elimination mechanism.
    \item Ammonia and amines react with acid chlorides to give amides, also through the addition–elimination mechanism of nucleophilic acyl substitution.
\end{enumerate}

SUMMARY Reactions of Carboxylic Acids

\begin{example} % 20-29
    Arrange each group of compounds in order of increasing basicity.
\end{example}

\begin{example} % 20-32
    Arrange each group of compounds in order of increasing acidity.
\end{example}

\begin{example} % 20-33
    What do the following $\mathrm{p}K_{\mathrm{a}}$ values tell you about the electron-withdrawing abilities of nitro, cyano, chloro, and hydroxyl  groups?
\end{example}

\begin{example} % 20-34
    Given the structure of ascorbic acid (vitamin C):
    \chemfig{
    *5(
        O                           % ① 环内氧（内酯的醚氧）
        -([::60]=O)                 % ② 羰基 =O（碳不标）
        =(-[::-60]OH)               % ③ 烯二醇一端 + —OH
        -(-[::-60]OH)               % ④ 烯二醇另一端 + —OH
        -(                          % ⑤ 环上碳，连接侧链（碳不标）
            [::60]                  % 侧链第1个碳（不标）
            -[::-90](<[::60]OH)      % 带 —OH（手性）
            -[::-60]OH              % 侧链末端 —CH₂OH（碳不标，只标 —OH）
        )-
    )
}
    \begin{enumerate}[label=(\alph*)]
        \item Is ascorbic acid a carboxylic acid?
        \item Compare the acid strength of ascorbic acid ($\mathrm{p}K_{\mathrm{a}} = 4.71$) with acetic acid. 
        \item Predict which proton in ascorbic acid is the most acidic. (lose proton and compare the resonance forms)
        \item Draw the form of ascorbic acid that is present in the body (aqueous solution, $\mathrm{pH} = 7.4$).
    \end{enumerate}
\end{example}

\begin{example} % 20-28 !!!
    Show how you would use extractions with a separatory funnel to separate a mixture of the following compounds: benzoic  acid, phenol, benzyl alcohol, aniline.

    sort first.
\end{example}

\begin{example} % 20-40
    The following NMR spectra correspond to compounds of formulas (A) (B) and (C)  respectively. Propose structures, and show how they are consistent with the observed absorptions.
\end{example}